\documentclass[final]{beamer}

\usepackage[scale=1.0]{beamerposter}
\usetheme{confposter}

\usepackage{graphicx}
\usepackage{booktabs}
\usepackage{hyperref}
\usepackage{amsmath}
\usepackage{multirow}

\setlength{\paperwidth}{36in}
\setlength{\paperheight}{48in}
\setlength{\topmargin}{-0.5in}
\usepackage[utf8]{inputenc}
\usepackage{amsmath}
\usepackage{amssymb}
\usepackage{amsfonts}
\usepackage{multicol}
\usepackage{listings}
\usepackage{xcolor}

\definecolor{lightgray}{gray}{0.95}

\lstset{
    basicstyle=\ttfamily\small,
    breaklines=true,
    columns=fullflexible,
    frame=single,
    backgroundcolor=\color{lightgray},
  % rulecolor=\color{black},
  % tabsize=4,
  showspaces=false,
   showstringspaces=false,
    language=c++,
     literate={'}{{'}}1
} 

\makeatletter
\renewcommand{\section}[1]{%
  \par\vspace{1em}%
  {\noindent\usebeamerfont{section title}\Large\bfseries #1\par}%
  \vspace{0.5em}\hrule\vspace{1em}%
}
\makeatletter

\makeatletter
\renewcommand{\subsection}[1]{%
  \par\vspace{1em}%
  {\noindent\usebeamerfont{subsection title}\large\bfseries #1\par}%
  \vspace{0.5em}\hrule\vspace{1em}%
}
\makeatletter




\title{C++ QUICK REFERENCE }
%\author{CS 3911}
\begin{document}

%\section*{C++ Reference Card}


\begin{frame}[t,fragile]{C++98}
%\begin{columns}[t,totalwidth=0.95\paperwidth, colsep=0.05\paperwidth]
\setlength{\columnseprule}{0.4pt}  % thickness of line

\begin{multicols}{4}[\setlength{\columnsep}{1.5cm}]

% ---------------- Column 1 ----------------
%\begin{column}{0.22\paperwidth}

%%begin{block}
\vspace{-1cm}
\section{Preprocessor}
\vspace{-1cm}
%begin{block}
\begin{lstlisting}[language=C++]
// Comment to end of line
/* Multi-line comment */
#include <stdio.h>         // Insert standard header file
#include "myfile.h"         // Insert file in current directory
#define X some text         // Replace X with some text
#define F(a,b) a+b          // Replace F(1,2) with 1+2
#define X \
 some text                 // Multiline definition
#undef X                    // Remove definition
#if defined(X)              // Conditional compilation (#ifdef X)
#else                       // Optional (#ifndef X or #if !defined(X))
#endif                      // Required after #if, #ifdef
\end{lstlisting}
%\end{block}

 \vspace{-1cm}
\section{Data Types}
\vspace{-1cm}
\begin{tabular}{|l|l|}
\hline
\textbf{Data Type} & \textbf{Description} \\
\hline
bool & boolean (true or false) \\
\hline
char & character ('a', 'b', etc.) \\
\hline
char[] & character array (C-style string  \\
 & if null terminated) \\
\hline
string & C++ string (from the STL) \\
\hline
int & Integer (1, 2, -1, 1000, etc.) \\
\hline
long int & long integer \\
\hline
float & single precision floating point \\
\hline
double & double precision floating point \\
\hline
\end{tabular}

These are the most commonly used types: this is not a complete list.
%%\end{block}

%begin{block}
 
\section{Operators}
The most commonly used operators in order of precedence:
\begin{enumerate}
    \item $++(\text{post-increment}), --(\text{post-decrement})$
    \item $!(\text{not}), ++ (\text{pre-increment}), -- (\text{pre-decrement})$
    \item $* / , \% (\text{modulus})$
    \item $+-$
    \item $<, <=, >, >=\text{}$
    \item $==(\text{equal-to}), != (\text{not-equal-to})$
    \item $\&\&$ (and)
    \item $||$ (or)
    \item $= (\text{assignment}), * , / , \%= , +=, -= $
\end{enumerate}
%\end{block}


\section{Looping}

%\begin{sectionblock}{Loops}
%\end{sectionblock}
%begin{block}
 
\subsection{while Loop}
\begin{lstlisting}
while (expression)
  statement;
\end{lstlisting}
{\it Example:}
\begin{lstlisting}
while (x < 100)
  cout << x++ << endl;
\end{lstlisting}
\begin{lstlisting}
while (expression){
  statement;
  statement;
}
\end{lstlisting}
{\it Example:}
\begin{lstlisting}
while (x < 100){
  cout << x << endl;
  x++;
}
\end{lstlisting}
%\end{block}

%begin{block}

\subsection{do-while Loop}
\begin{lstlisting}
do
  statement;
while (expression);
\end{lstlisting}
{\it Example:}
\begin{lstlisting}
do
  cout << x++ << endl;
while (x < 100);
\end{lstlisting}
\begin{lstlisting}
do{
  statement;
  statement;
} while (expression);
\end{lstlisting}
{\it Example:}
\begin{lstlisting}
do {
  cout << x << endl;
  x++;
} while (x < 100);
\end{lstlisting}
%\end{block}

%\end{column}

% ---------------- Column 2 ----------------
%\begin{column}{0.3\paperwidth}
\subsection{For Loop}
%begin{block}

%\begin{tabular}{l | l}

%\begin{minipage}{0.45\linewidth}

\begin{lstlisting}
for (initialization; test; update)
  statement;
\end{lstlisting}

\begin{lstlisting}
for (initialization; test; update){
  statement;
  statement;
}
\end{lstlisting}


%\end{minipage}
%&
%\begin{minipage}{0.45\linewidth}

%\end{column}
% ---------------- Column 2 ----------------
%\begin{column}{0.23\paperwidth}

{\it Example:}

\begin{lstlisting}
for (count = 0; count < 10; count++) {
  cout << "count equals: " << count << endl;
  
}
\end{lstlisting}

%\end{minipage}

%\end{tabular}
%\end{block}



%\end{column}
% ---------------- Column 2 ----------------
%\begin{column}{0.3\paperwidth}

%%begin{block}
\vspace{-1.5cm}
 \subsection{ Decision Statements}
%%\end{block}
\vspace{-1.5cm}



%begin{block}
 
\subsection{if}
\begin{lstlisting}
if (expression)
  statement;
\end{lstlisting}
Example:
\begin{lstlisting}
if (x < y)
  cout << x;
\end{lstlisting}
%\end{block}




%begin{block}
 
\subsection{if/else}
\begin{lstlisting}
if (expression)
  statement;
else
  statement;
\end{lstlisting}
Example:
\begin{lstlisting}
if (x < y)
  cout << x;
else
  cout << y;
\end{lstlisting}
%\end{block}

%begin{block}
 
\subsection{switch/case}
\begin{lstlisting}
switch(int expression)
{
  case int-constant:
    statement(s);
    break;
  case int-constant:
    statement(s);
    break;
  default:
    statement;
}
\end{lstlisting}
{\it Example:}
\begin{lstlisting}
switch(choice)
{ 
  case 0:
    cout << "Zero";
    break;
  case 1:
    cout << "One";
    break;
  default:
    cout << "What?";
}
\end{lstlisting}
%\end{block}

\section{Pointers}
A pointer variable (or just pointer) is a variable that stores a memory address. Pointers allow the indirect manipulation of data stored in memory. Pointers are declared using *. To set a pointer's value to the address of another variable, use the \& operator.
Example:
\begin{lstlisting}
char c = 'a';
char* cPtr;
cPtr = &c;
\end{lstlisting}
Use the indirection operator (*) to access or change the value that the pointer references.
Example:
\begin{lstlisting}
// continued from example above
*cPtr = 'b';
cout << *cPtr << endl; // prints the char b
// prints the char b
\end{lstlisting}

Array names can be used as constant pointers, and pointers can be used as array names.
Example:
\begin{lstlisting}
int numbers[] = {10, 20, 30, 40, 50};
int* numPtr = numbers;
cout << numbers[0] << endl; // prints 10
cout << *numPtr << endl;    // prints 10
cout << numbers[1] << endl; // prints 20
cout << *(numPtr + 1) << endl; // prints 20
cout << numPtr[2] << endl;  // prints 30
\end{lstlisting}
%%\end{block}
\vspace{-1.5cm}
\section{Dynamic Memory}
\vspace{-1.5cm}

%begin{block}
 \subsection{ Allocate Memory}
\begin{tabular}{|l|l|}
\hline
\textbf{Syntax} & \textbf{Examples} \\
\hline
ptr = new type; & int* iPtr; \\
 & iPtr = new int; \\
\hline
ptr = new type [size]; & int* intArray; \\
 & intArray = new int[5]; \\
\hline
\end{tabular}
%\end{block}

%begin{block}
 %\end{column}
%\begin{column}{0.23\paperwidth}

\subsection{Deallocate Memory}
\begin{tabular}{|l|l|}
\hline
\textbf{Syntax} & \textbf{Examples} \\
\hline
delete ptr; & delete iPtr; \\
\hline
delete [] ptr; & delete [] intArray; \\
\hline
\end{tabular}

Once a pointer is used to allocate the memory for an array, array notation can be used to access the array locations.
Example:
\begin{lstlisting}
int* intArray;
intArray = new int[5];
intArray[0] = 23;
intArray[1] = 32;
\end{lstlisting}
%\end{block}
\section{I/O}
%begin{block}
 \subsection{ Console Input/Output}
\begin{itemize}
    \item \texttt{cout <<} console out, printing to screen
    \item \texttt{cin >>} console in, reading from keyboard
    \item \texttt{cerr <<} console error
\end{itemize}
Example:
\begin{lstlisting}
cout << "Enter an integer: ";
cin >> i;
cout << "Input: " << i << endl;
\end{lstlisting}
%\end{block}

%begin{block}
 \subsection{ File Input/Output}
Example (input):
\begin{lstlisting}
ifstream inputFile;
inputFile.open("data.txt");
inputFile >> inputVariable;
// you can also use get (char) or
// getline (entire line) in addition to >>
inputFile.close();
\end{lstlisting}
Example (output):
\begin{lstlisting}
ofstream outFile;
outFile.open("output.txt");
outFile << outputVariable;
outFile.close();
\end{lstlisting}
%\end{block}
\section{Function}
%begin{block}
 \subsection{Functions}
Functions return at most one value. A function that does not return a value has a return type of \texttt{void}. Values needed by a function are called parameters.
\begin{lstlisting}
return_type function(type p1, type p2,...)
{
  statement;
  statement;
}
\end{lstlisting}
Examples:
\begin{lstlisting}
int timesTwo(int v) {
  int d;
  d = v * 2;
  return d;
}
void printCourseNumber() {
  cout << "CSE1284" << endl;
  return;
}
\end{lstlisting}
%\end{block}

%begin{block}
 \subsection{Passing Parameters by Value}
\begin{lstlisting}
return_type function(type p1)
\end{lstlisting}
Variable is passed into the function but changes to \texttt{p1} are not passed back.
%\end{block}

%begin{block}
 \subsection{Passing Parameters by Reference}
\begin{lstlisting}
return_type function(type &p1)
\end{lstlisting}
Variable is passed into the function and changes to \texttt{p1} are passed back.
%\end{block}


%begin{block}
 \subsection{Default Parameter Values}
\begin{lstlisting}
return_type function(type p1=val)
\end{lstlisting}
\texttt{val} is used as the value of \texttt{p1} if the function is called without a parameter.
%\end{block}
%\end{column}
% ---------------- Column 3 ----------------
%\begin{column}{0.23\paperwidth}
%begin{block}
\vspace{-1.cm}
 \subsection{Struct/Classes}
 \vspace{-1.cm}
%\end{block}
%begin{block}
%\end{block}

%begin{block}
 \subsection{Declaration}
\begin{lstlisting}
class classname {
public:
  classname (params);
  ~classname();
  type member1;
  type member2;
protected:
  type member3;
private:
  type member4;
};
\end{lstlisting}
Example:
\begin{lstlisting}
class Square {
public:
  Square();
  Square (float w);
  void setWidth(float w);
  float getArea();
private:
  float width;
};
\end{lstlisting}
\begin{itemize}
    \item \texttt{public} members are accessible from anywhere the class is visible.
    \item \texttt{private} members are only accessible from the same class or a friend (function or class).
    \item \texttt{protected} members are accessible from the same class, derived classes, or a friend (function or class).
    \item constructors may be overloaded just like any other function. You can define two or more constructors as long as each constructor has a different parameter list.
\end{itemize}
%\end{block}

%begin{block}
 \subsection{Definition of Member Functions}
\begin{lstlisting}
return_type classname::functionName (params) {
  statements;
}
\end{lstlisting}
Examples:
\begin{lstlisting}
Square::Square() {
  width = 0;
}
void Square::setWidth(float w) {
  if (w >= 0)
    width = w;
  else
    exit(-1);
}
float Square::getArea() {
  return width * width;
}
\end{lstlisting}
%\end{block}

%begin{block}
 \subsection{Definition of Instances}
\begin{lstlisting}
classname varName;
classname* ptrName;
\end{lstlisting}
Example:
\begin{lstlisting}
Square s1();
Square s2(3.5);
Square* sPtr;
sPtr = new Square (1.8);
\end{lstlisting}
%\end{block}

%begin{block}
 \subsection{Accessing Members}
\begin{lstlisting}
varName.member = val;
varName.member();
ptrName->member = val;
ptrName->member();
\end{lstlisting}
Example:
\begin{lstlisting}
s1.setWidth(1.5);
cout << s1.getArea();
cout << sPtr->getArea();
\end{lstlisting}
%\end{block}
%\end{column}
%\end{columns}
\end{multicols}
\end{frame}
\begin{frame}[t,fragile]{OOP, Templates, Exceptions}

\setlength{\columnseprule}{0.4pt}  % thickness of line

\begin{multicols}{4}[\setlength{\columnsep}{1.5cm}]

%\begin{columns}[t,totalwidth=0.95\paperwidth, colsep=0.05\paperwidth]

% ---------------- Column 1 ----------------
%\begin{column}{0.22\paperwidth}
%begin{block}

 \subsection{OOP Principles}
\textbf{Encapsulation:} Hide internal implementation details by using access specifiers (\texttt{public}, \texttt{private}, \texttt{protected}) to control access to member variables and functions.

\textbf{Abstraction:} Define interfaces that represent only essential features while hiding complexity. Use abstract classes with pure virtual functions.
%\end{block}

%begin{block}

 \subsection{Inheritance}
Inheritance allows a new class to be based on an existing class. The new class inherits all the member variables and functions (except the constructors and destructor) of the class it is based on.

Example:
\begin{lstlisting}
class Student
{
public:
  Student (string n, string id);
  void print();
protected:
  string name;
  string netID;
};
class GradStudent : public Student
{
public:
  GradStudent(string n, string id, string prev);
  void print();
protected:
  string prevDegree;
};
\end{lstlisting}
%\end{block}

%begin{block}
 \subsection{Visibility of Members after Inheritance}
\begin{tabular}{|l|l|l|l|}
\hline
\multirow{2}{*}{Inheritance Specification} & \multicolumn{3}{c|}{Access Specifier in Base Class} \\
\cline{2-4}
 & private & protected & public \\
\hline
private & private & private & private \\
\hline
protected & protected & protected & protected \\
\hline
public & protected & protected & public \\
\hline
\end{tabular}
%\end{block}

%begin{block}
 \subsection{Polymorphism}
\textbf{Compile-time (Static) Polymorphism:} Function overloading and templates.

\textbf{Runtime (Dynamic) Polymorphism:} Use virtual functions to allow derived classes to override base class methods. Write code that works with base class pointers/references to handle derived classes dynamically.

\textbf{Example:}
\begin{lstlisting}
class Shape // Abstract class
{
public:
  virtual void draw() = 0; // pure virtual
  virtual double area() = 0;
  virtual ~Shape() {}
};

class Circle : public Shape
{
public:
  void draw() { /* implementation */ }
  double area() { return 3.14 * r * r; }
private:
  double r;
};

Shape* shapes[10];
shapes[0] = new Circle();
shapes[0]->draw(); // calls Circle::draw()
\end{lstlisting}
%\end{block}

%begin{block}
 \subsection{Operator Overloading}
C++ allows you to define how standard operators (+, -, *, etc.) work with classes that you write. For example, to use the operator $+$ with your class, you would write a function named \texttt{operator+} for your class.
Example:
Prototype for a function that overloads $+$ for the \texttt{Square} class:
\begin{lstlisting}
Square operator+ (const Square &);
\end{lstlisting}

%\end{column}
% ---------------- Column 2 ----------------
%\begin{column}{0.23\paperwidth}
If the object that receives the function call is not an instance of a class that you wrote, write the function as a \texttt{friend} of your class. This is standard practice for overloading \texttt{<<} and \texttt{>>}.






Example:
Prototype for a function that overloads \texttt{<<} for the \texttt{Square} class:
\begin{lstlisting}
friend ostream & operator<<
  (ostream &, const Square &);
\end{lstlisting}
Make sure the return type of the overloaded function matches what C++ programmers expect. The return type of relational operators $(<, >, ==, \text{etc.})$ should be \texttt{bool}, the return type of \texttt{<<} should be \texttt{ostream \&}, etc.
%\end{block}


\subsection{Namespaces}
\begin{lstlisting}[language=C++]
namespace N { class T {}; }
N::T t;
using namespace N;
\end{lstlisting}


%begin{block}
 \subsection{Exceptions}
Example:
\begin{lstlisting}
try
{
  // code here calls functions that might
  // throw exceptions
  quotient = divide (num1, num2);

  // or this code might test and throw
  // exceptions directly
  if (num3 < 0)
    throw -1; // exception to be thrown can
              // be a value or an object
}
catch (int)
{
  cout << "num3 can not be negative!";
  exit(-1);
}
catch (char* exceptionString)
{
  cout << exceptionString;
  exit(-2);
}
// add more catch blocks as needed
\end{lstlisting}
%\end{block}

%begin{block}
 \subsection{Function Templates}
Example:
\begin{lstlisting}
template <class T>
T getMax(T a, T b)
{
  if (a > b)
    return a;
  else
    return b;
}
\end{lstlisting}
Example calls:
\begin{lstlisting}
int a = 9, b = 2, c;
c = getMax(a, b);
float f = 5.3, g = 9.7, h;
h = getMax(f, g);
\end{lstlisting}
%\end{block}

%begin{block}
 \subsection{Class Templates}
Example:
\begin{lstlisting}
template <class T>
class Point
{
public:
  Point(T x, T y);
  void print();
  double distance (Point<T> p);
private:
  T x;
  T y;
};
\end{lstlisting}
Examples:
\begin{lstlisting}
Point<int> p1(3, 2);
Point<float> p2(3.5, 2.5);
p1.print();
p2.print();
\end{lstlisting}
%\end{block}

%begin{block}
 \subsection{Suggested Websites}
\begin{itemize}
    \item C++ Reference: \texttt{http://www.cppreference.com/}
    \item C++ Tutorial: \texttt{http://www.cplusplus.com/doc/tutorial/}
\end{itemize}
%Developed for Mississippi State University's CSE1284 and CSE1384 courses.
%Download the latest version from \texttt{http://cse.msstate.edu/~crumpton/reference}.  
%February 17, 2009.
%\end{block}

%\end{column}

%\end{columns}
\end{multicols}
\end{frame}

\end{document}